\documentclass{article}

\title{Chem 001A\\ Week 8}
\author{Danny Topete}
\date{November 17, 2025}

\begin{document}
\maketitle

\section{Monday Lecture}

\subsection{Determining molecular formulas}

\begin{itemize}
  \item \textbf{Molecular formula:} provides absolute number of atoms of each element in a molecule
  \item \textbf{Empirical Formula:} provides relative number of atoms of each element in a molecule
  \item You need the empirical formula in order to build up the molecular formula
  \item In the example of Benzine, you need 6 empiricla formulas in order to achieve the molecular formula ($C_6H_6$)
\end{itemize}

\subsection{Going over example}
\begin{itemize}
\item \textbf{Example:}
  $$ n = \frac{molarMass}{EmpiricalUnitMass CH_2O} = \frac{180.0 amu}{30.0 amu} = 6 $$
\item Carbohydrates contain $CH_2O $
\item Nicontine sample:

  C = 74.02\%, H = 8.710\%, N = 17.27\%. If 40.57g of Nicontine contains 0.2500 mol nicotine, what is the
  the molecular form.

  72.02g C $\frac{1 mol C}{ 12.01g gC} = 6.163 mol C$

  Keep going, for the mols for each atom, then divide by smallest. You get the proportions; an empirical formula.
\end{itemize}

\subsection{Introduction to Molarity}
\begin{itemize}
  \item Being able to quantify solutions
  \item \textbf{Molarity:} Is defined as numbers of moles of solute in exactly 1 liter of solution
    $$ M = \frac{mol Solute}{1L solute} $$
  \item \textbf{Example:} 355ml soft drink, contains 0.133 mol concentration of sucrose.
    What is the molar concentration of sucrose in the beverage?
    (I assume they looked at the amount in liters, then multiplies by mols of sugar / Listers)
  \item Molarity really is just a definition, and they are applying dimensional analysis upon this.
  \item Do lots of practice of making solutions with molarity
\end{itemize}

\subsection{Preparing Solutions}
\begin{itemize}
  \item Builds off molarity, you need to know
    how many moles are required to build this solution
\end{itemize}

\end{document}


